%% CAPITULO 2 -- REFERENCIAL TEORICO (esqueleto da dissertação)

\section{Referencial teórico}

A avaliação de sistemas RAG tem se apoiado em métricas automáticas que separam
a qualidade da recuperação da qualidade da geração. Ao propor o RAGAS, um
\textit{framework} de avaliação sem referência humana, \citeonline{es2023ragas}
resumem por que essa avaliação é difícil:

\begin{citacao}
Evaluating RAG architectures is, however, challenging because there are several
dimensions to consider: the ability of the retrieval system to identify relevant
and focused context passages, the ability of the LLM to exploit such passages in
a faithful way, or the quality of the generation itself.
\parencite{es2023ragas}
\end{citacao}

As dimensões listadas pelos autores --- recuperação, fidelidade e qualidade da
geração --- não incluem a capacidade de a resposta indicar o documento, a unidade
documental e a vigência da norma que a sustenta. No setor público essa lacuna
pesa porque, como define \citeonline{bovens2007accountability}, a
\textit{accountability} é uma relação em que o ator tem a obrigação de explicar
e justificar sua conduta perante um fórum que pode questioná-lo e julgá-lo; uma
resposta sobre documento público que não permite localizar a sua evidência não
sustenta essa obrigação. A Tabela~\ref{tab:dimensoes} reúne as cinco dimensões
do modelo proposto, que acrescentam rastreabilidade e adequação institucional às
dimensões técnicas.

\begin{table}[htb]
  \caption{Dimensões do modelo multidimensional de avaliação}
  \label{tab:dimensoes}
  \centering
  \begin{tabular}{p{4.2cm}p{10.4cm}}
    \hline
    \textbf{Dimensão} & \textbf{O que observa} \\
    \hline
    Recuperação & relevância e cobertura dos trechos recuperados \\
    Geração & correção e fidelidade da resposta ao contexto recuperado \\
    Eficiência & custo em \textit{tokens}, tempo de indexação e latência \\
    Rastreabilidade & IRD: documento, unidade documental, metadados,
      auditabilidade, vigência e não extrapolação \\
    Adequação institucional & abstenção diante de evidência insuficiente e
      requisitos do regime público \\
    \hline
  \end{tabular}
  \fonte{Elaborada pelo autor (2026).}
\end{table}
